\chapter{Incision and Drainage}

\section*{Overview}
Incision and drainage (I\&D) is a minor surgical procedure to release pus from an abscess, boil, or infected cyst. It relieves pressure, reduces pain, and allows the infection to heal.

\section*{Alternative Names}
I\&D, Abscess drainage

\section*{Principle}
The skin and subcutaneous tissue over an abscess are incised to open the cavity, pus is drained, and loculations are broken down. The wound is often left open or packed so infection can drain while the cavity granulates and heals from the base outward.
\section*{History}
The surgical maxim 'Ubi pus, ibi evacua' (Where there is pus, evacuate it) has been followed since the Roman physician Celsus in the 1st century AD.

\section*{Why This Test or Procedure Is Ordered}
Ordered for fluctuant skin abscesses or localized purulent infections that do not resolve with antibiotics alone.

\section*{Is This a Primary or Secondary Test or Procedure}
Primary definitive treatment for localized, fluctuant soft-tissue abscesses.

\section*{How This Test or Procedure Is Performed}
The area is cleaned, and local anesthetic is injected. An incision is made over the center of the abscess. The pus is drained, loculations are broken, and the cavity may be packed with gauze.

\section*{What Preparation Is Needed}
Irrigate the site. Obtain scalpel, packing gauze, forceps, and local anesthetic.

\section*{Contraindications and Precautions}
Facial abscesses in the 'danger triangle' (risk of cavernous sinus thrombosis) or abscesses close to major blood vessels may require surgical referral.

\section*{Risks}
Bleeding, pain, recurrence, scar formation, or spread of infection.

\section*{What Happens After the Test or Procedure}
Keep the wound clean and covered. If packed, the packing is removed or changed in 24-48 hours. Warm compresses are encouraged.

\section*{Understanding Your Results}
Complete drainage leads to immediate pain relief, decreased swelling, and resolution of local infection.

\section*{Normal Results}
Evacuation of purulent fluid; collapse of the abscess cavity; immediate relief of pressure pain.

\section*{Follow-up Tests and Procedures}
Wound check and packing removal in 24-48 hours; monitoring for systemic signs of infection (fever).

\section*{References}
\begin{enumerate}
  \item MedlinePlus. Incision and Drainage. U.S. National Library of Medicine; 2025.
  \item Current Medical Diagnosis and Treatment. McGraw-Hill; 2025.
\end{enumerate}

\clearpage
